\label{sec:newstates}

\subsection{Definition of the New-States Paradox}

\begin{definition}[New-States Paradox]
  An apportionment method $M$ satisfies the New-States Paradox if
  adding a new state $k$ with population $p_k$ and
  $m = \text{round}(p_k \times H / N)$ additional seats
  (so the new House size is $H' = H + m$) causes existing states
  to exchange seats: state $i$ gains a seat and state $j$ loses one,
  even though no change was made to states $i$ and $j$ themselves.
\end{definition}

The New-States Paradox is paradoxical because the new state's seats are
supposed to be ``proportional'' to its population ---
the existing states should not be affected.
But under Hamilton, the addition of the new state changes the exact quotas
of all states (because the national total $N$ and divisor change),
and the remainder-seat ranking can shift.

\subsection{The Oklahoma 1907 Example}

When Oklahoma was admitted to the Union in 1907, Congress increased
the House size from 386 to 391 (by 5 seats, approximately proportional
to Oklahoma's population).
Under Hamilton:

\begin{enumerate}
  \item Before Oklahoma: $H = 386$, $N \approx 76{,}000{,}000$.
    Maine received 3 seats; New York received 38 seats.
  \item After Oklahoma: $H = 391$, $N \approx 77{,}000{,}000$.
    The new Hamilton apportionment (with Oklahoma receiving 5 seats)
    gave Maine 4 seats and New York 37 seats.
\end{enumerate}

The addition of Oklahoma with its proportional 5 seats caused Maine
to gain a seat and New York to lose a seat, despite no change in
Maine's or New York's populations.
This is the New-States Paradox \citep{balinski1982fair}.

\subsection{Numerical Example}

\begin{example}[3-State New-States Paradox]
  Start with 2 states: $\vec{p} = (30, 20)$, $N = 50$, $H = 5$.
  Quotas: $(3.0, 2.0)$, lower quotas $(3, 2)$, $r = 0$.
  Hamilton: $(3, 2)$.

  Add state 3 with $p_3 = 10$.
  Proportional seats: $10 \times 5 / 50 = 1$.
  New House size: $H' = 6$; new total population: $N' = 60$.

  New quotas: $(30 \times 6/60, 20 \times 6/60, 10 \times 6/60)
  = (3.0, 2.0, 1.0)$.
  Lower quotas: $(3, 2, 1)$, $r = 0$.
  Hamilton: $(3, 2, 1)$.
  No paradox in this example.

  For a paradox, we need the remainder-seat ranking to change.
  Consider $\vec{p} = (34, 66)$, $N = 100$, $H = 10$.
  Quotas: $(3.4, 6.6)$, lower quotas $(3, 6)$, $r = 1$.
  Fractional parts: $(0.4, 0.6)$; state 2 gets remainder.
  Hamilton: $(3, 7)$.

  Add state 3 with $p_3 = 7$. Proportional seats: $7 \times 10/100 = 0.7 \approx 1$.
  $H' = 11$, $N' = 107$.
  New quotas: $(34 \times 11/107, 66 \times 11/107, 7 \times 11/107)
  = (3.495, 6.785, 0.720)$.
  Lower quotas: $(3, 6, 0)$, sum = 9, $r = 2$.
  Fractional parts: $(0.495, 0.785, 0.720)$.
  State 2 gets first remainder (0.785), state 3 gets second (0.720).
  Hamilton: $(3, 7, 1)$.
  No change for states 1 and 2.

  Adjusting population so that state 1's fractional part displaces state 2:
  populations $(34, 66)$ with $H = 10$: state 1 has fractional 0.4,
  state 2 has 0.6.
  Add state 3 with $p_3 = 50$: proportional $= 50 \times 10/100 = 5$ seats.
  $H' = 15$, $N' = 150$.
  New quotas: $(34\times 15/150, 66\times 15/150, 50\times 15/150)
  = (3.4, 6.6, 5.0)$.
  Same fractions as before! No paradox.

  The New-States Paradox requires specific population configurations;
  the Oklahoma example uses the actual 1907 census data and is
  documented in \citet{balinski1982fair}.
\end{example}

\subsection{Why All Three Paradoxes Share the Same Root}

The Alabama, Population, and New-States paradoxes all arise from
the same mathematical structure: Hamilton's fractional-part ranking
is not monotone in the parameters (House size, populations, number
of states) that are supposed to determine the apportionment.
Specifically:
\begin{itemize}
  \item The fractional parts $f_i = q_i - \lfloor q_i \rfloor$
    are discontinuous functions of $H$ and $\vec{p}$
    (a small change in $H$ or $p_i$ can cross an integer boundary,
    dramatically changing $f_i$).
  \item The ranking of fractional parts can change even when no state's
    fractional part changes monotonically.
  \item States that should benefit from an increase in $H$ or $p_i$
    can instead lose their remainder seat because another state's
    fractional part happened to increase more.
\end{itemize}
Divisor methods avoid all of this because they use a priority function
$P_i(s)$ that is monotone in $p_i$ and the seat allocation is determined
by a priority queue that only grows (never reconsidering previously
allocated seats).
